\section*{Chapitre \ref{ch:m3:power-markets-trading} --- Marchés de l'électricité II~: négociation}

\begin{solution}{exo:m3:power-markets-trading:1}
\omterm{def:m3:power-markets-trading:capture}{Prix capté} $(800 + 800 + 400 + 900)/100 = \qty{29.00}{EUR/MWh}$~; charge de base
\qty{50.00}{EUR/MWh}~; taux de capture 58~\%.
\end{solution}

\begin{solution}{exo:m3:power-markets-trading:2}
Il verse $20 \times 300 = $ 6\,000~EUR. En position longue de 20~MWh, il aurait reçu 6\,000~EUR sous un
prix unique.
\end{solution}

\begin{solution}{exo:m3:power-markets-trading:3}
Les positions sont arrêtées la veille sur des prévisions qui ne cessent de changer~: vent, solaire, demande et
indisponibilités de centrales. Le trading infrajournalier permet à chaque partie de corriger sa position à mesure que sa prévision s'améliore,
au lieu de payer des \omterm{def:m3:power-markets-trading:imbalance}{prix de règlement des écarts} sur l'ensemble de l'erreur.
\end{solution}

\begin{solution}{exo:m3:power-markets-trading:4}
$50 \times (15 - (-5)) \times 10 = $ 10\,000~EUR.
\end{solution}

\begin{solution}{exo:m3:power-markets-trading:5}
Elle prélève $200/\sqrt{0.88} = 213.2$~MWh à \qty{20}{}, en payant 4\,264~EUR, et restitue
$200\sqrt{0.88} = 187.6$~MWh à \qty{120}{}, en recevant 22\,514~EUR~: soit 18\,250~EUR.
\end{solution}

\begin{solution}{exo:m3:power-markets-trading:6}
$(60 - 46.07) \times 100\,000 = $ 1.39~million~EUR par an.
\end{solution}

\begin{solution}{exo:m3:power-markets-trading:7}
65\,968~EUR par MW en 2024 et 74\,546~EUR en 2025 avec un cycle par jour~; 88\,688~EUR en 2025
avec deux.
\end{solution}

\begin{solution}{exo:m3:power-markets-trading:8}
Le backtest suppose une connaissance parfaite des prix du jour, une batterie preneuse de prix qui
n'influe pas sur l'enchère, et suppose l'absence de dégradation, d'indisponibilité ou de frais~; les opérateurs réels prévoient, et leurs
propres offres réduisent les écarts qu'ils négocient. Il ignore aussi les revenus infrajournaliers et de réserve qui
peuvent peser davantage. Il faut le considérer comme une borne supérieure des revenus sur le \omterm{def:m3:power-markets-design:da}{marché J-1}.
\end{solution}

\begin{solution}{pb:m3:power-markets-trading:1}
\textbf{1.} Elle se charge de 14~h à 16~h (prix $-2.30$ et $-2.17$) et se décharge de 20~h
à 22~h (140.99 et 175.32)~: la paire d'heures la moins chère suivie de la plus chère.
\textbf{2.} Elle prélève 213.2~MWh et restitue 187.6~MWh. \textbf{3.} 30\,149~EUR. \textbf{4.} À un
prix négatif, elle est payée pour charger~; les prix les plus bas étaient négatifs cet après-midi.
\textbf{5.} 760~EUR de plus~: le deuxième meilleur écart de la journée était faible.
\textbf{6.} 65\,968~EUR en 2024 et 74\,546~EUR en 2025 par MW. \textbf{7.} Davantage d'heures négatives et
proches de zéro en milieu de journée à mesure que le solaire s'est développé, et des prix élevés le soir~: l'écart quotidien s'est élargi.
\textbf{8.} 14\,142~EUR par MW. \textbf{9.} Les écarts sont les plus larges lorsque le solaire est fort et que les soirées
sont tendues (printemps et été) et lors des semaines d'hiver volatiles. \textbf{10.} Sans pertes, elle aurait
gagné 32\,078~EUR le 18 juin~: la perte de 12~\% coûte environ 1\,929~EUR.
\textbf{11.} Les programmes réels se construisent sur des prévisions~; les offres d'un faiseur de prix déplacent aussi les prix qu'il
négocie. \textbf{12.} Arbitrage infrajournalier, réserves d'équilibrage (rémunérations de capacité), optimisation des
écarts pour un portefeuille et services réseau. \textbf{13.} Elles se resserrent~: davantage d'acheteurs à midi
relèvent les prix bas, davantage de vendeurs le soir abaissent les prix élevés. \textbf{14.} Chaque cycle
use les cellules~; limiter les cycles revient à arbitrer entre le revenu d'aujourd'hui et la capacité de demain. \textbf{15.} Davantage
de pas de prix plus courts à arbitrer, et une meilleure adéquation à la réponse rapide de la batterie.
\textbf{16.} Les deux résultent de la même chute des prix en milieu de journée~: ce qui abaisse le \omterm{def:m3:power-markets-trading:capture}{prix capté} du solaire
élargit l'écart sur lequel négocient les batteries. \textbf{17.} Autres batteries, stations de transfert d'énergie par pompage, demande industrielle flexible,
électrolyseurs et exportateurs. \textbf{18.} Seulement tant que le solaire continue de dépasser la flexibilité~;
elle s'érode à mesure que le stockage se développe. \textbf{19.} \emph{Résultat à retenir~:} 74\,546~EUR par MW en 2025
(65\,968~EUR en 2024), un cycle par jour, prévoyance parfaite. \textbf{20.} Déplacer l'énergie dans le temps~:
elle achète les heures que d'autres ne peuvent pas utiliser et vend celles qu'ils ne peuvent pas servir.
\end{solution}

\begin{solution}{iq:m3:power-markets-trading:1}
Le prix auquel le gestionnaire du réseau règle l'écart d'une partie à son programme, dérivé de
l'énergie d'équilibrage qu'il a activée. Être en déficit coûte peu lorsque le système est en excédent~: le \omterm{def:m3:power-markets-trading:imbalance}{prix de règlement des
écarts} est alors bas (voire négatif), de sorte qu'acheter à l'opérateur l'énergie manquante coûte peu.
\iqlookfor{la direction du système fixe le prix~; le signe de l'erreur par rapport à celui du système.}
\end{solution}

\begin{solution}{iq:m3:power-markets-trading:2}
Le solaire produit aux mêmes heures que tout le reste du solaire, et les prix de ces heures baissent à mesure que la production solaire
augmente~; le prix pondéré par la production est donc inférieur à la moyenne temporelle. Chaque nouveau parc solaire approfondit
la baisse des prix de midi, si bien que le taux de capture continue de baisser à moins que le stockage ou la demande ne l'absorbe.
\iqlookfor{autocannibalisation.}
\end{solution}

\begin{solution}{iq:m3:power-markets-trading:3}
Prévoir le profil horaire de production (météo, dégradation, écrêtement) et la courbe horaire de prix
sur dix ans (modèles fondamentaux ou de scénarios)~; valoriser la production à son \omterm{def:m3:power-markets-trading:capture}{prix capté}, et non
au prix de base~; ajouter les coûts d'écart et de profil~; actualiser~; comparer au prix du \omterm{def:m3:power-markets-trading:ppa}{contrat d'achat d'électricité} (PPA)~; tester
les scénarios de prix et la corrélation entre production et prix.
\iqlookfor{\omterm{def:m3:power-markets-trading:capture}{prix capté}, profil et corrélation, scénarios.}
\end{solution}

\begin{solution}{iq:m3:power-markets-trading:4}
Une position financière sur le \omterm{def:m3:power-markets-design:da}{marché J-1} à un nœud, réglée au prix temps réel~: elle arbitre
l'écart entre le \omterm{def:m3:power-markets-design:da}{marché J-1} et le temps réel, et aligne les deux. Elle peut servir à déplacer la congestion du \omterm{def:m3:power-markets-design:da}{marché J-1} pour
profiter aux FTR d'un trader, en perdant sur l'\omterm{def:m3:power-markets-trading:virtual}{offre virtuelle} et en gagnant davantage sur les droits~; les règles retiennent les
paiements de FTR lorsque ce schéma apparaît.
\iqlookfor{rôle de convergence, et manipulation croisée des FTR et du virtuel.}
\end{solution}

\begin{solution}{iq:m3:power-markets-trading:5}
La taille de l'erreur et, surtout, sa corrélation avec l'écart du système, donc avec le
\omterm{def:m3:power-markets-trading:imbalance}{prix de règlement des écarts}~: une erreur dans le sens du système coûte cher, une erreur à contre-courant est rémunérée. Un
portefeuille dont les erreurs sont corrélées à celles de la zone (toutes les prévisions éoliennes fausses ensemble) paie le plus.
\iqlookfor{corrélation avec le système, et pas seulement la variance.}
\end{solution}

\begin{solution}{iq:m3:power-markets-trading:6}
Une optimisation stochastique sur des scénarios de prix pour les trois marchés, sous contraintes de puissance,
d'énergie, de rendement et de cycles, des engagements de capacité de réserve qui réduisent la marge de manœuvre, et
des décisions séquentielles (\omterm{def:m3:power-markets-design:da}{marché J-1} d'abord, infrajournalier et réserves ensuite)~: un modèle à deux étapes ou
à horizon glissant, résolu sous forme de programme linéaire par ensemble de scénarios, avec un coût de dégradation par cycle.
\iqlookfor{co-optimisation, marchés séquentiels, contraintes et incertitude.}
\end{solution}
